\section{Method}
\label{sec:method}

\subsection{Problem Setting}

A knowledge graph $\mathcal{G} = (\mathcal{E}, \mathcal{R}, \mathcal{T})$ consists of entities $\mathcal{E}$, relations $\mathcal{R}$, and observed triples $\mathcal{T} \subseteq \mathcal{E} \times \mathcal{R} \times \mathcal{E}$. Given a query $(h, r, ?)$, link prediction ranks candidate tail entities by plausibility.

\paragraph{Coverage.}
We define the coverage function $\cov: \mathcal{E} \times \mathcal{R} \to \{0, 1\}$:
\[
\cov(e, r) = \begin{cases}
1 & \text{if } \exists t: (e, r, t) \in \mathcal{T} \text{ or } (t, r, e) \in \mathcal{T} \\
0 & \text{otherwise}
\end{cases}
\]
This indicates whether entity $e$ has been observed with relation $r$ in training.

\paragraph{Novel Context.}
A test triple $(h, r, t)$ is in \textbf{novel context} if $\cov(h, r) = 0$ or $\cov(t, r) = 0$. The entity is familiar (seen in training), but appears with a relation it has never been observed with.

\subsection{RCUE: Relation-Conditioned Uncertainty Estimation}

RCUE models uncertainty as a function of both entity and relation:

\paragraph{Entity and Relation Embeddings.}
Each entity $e \in \mathcal{E}$ has embedding $\mathbf{e} \in \mathbb{R}^d$. Each relation $r \in \mathcal{R}$ has embedding $\mathbf{r} \in \mathbb{R}^d$. These are used for both scoring and uncertainty estimation.

\paragraph{Scoring Function.}
We use DistMult scoring:
\[
s(h, r, t) = \langle \mathbf{h}, \mathbf{r}, \mathbf{t} \rangle = \sum_{i=1}^d h_i \cdot r_i \cdot t_i
\]
Other scoring functions (TransE, ComplEx) are compatible with our uncertainty estimation.

\paragraph{Relation-Conditioned Variance.}
The core innovation is computing entity variance conditioned on the relation:
\[
\sigma^2(e | r) = f_\theta([\mathbf{e}; \mathbf{r}]) \cdot \text{boost}(\cov(e, r))
\]

where $f_\theta$ is a neural network:
\[
f_\theta(\mathbf{x}) = \text{Softplus}(\text{MLP}(\mathbf{x}))
\]

The MLP has two hidden layers with ReLU activations. Softplus ensures positive variance.

\paragraph{Coverage Boost.}
The multiplicative boost amplifies uncertainty for unseen entity-relation pairs:
\[
\text{boost}(c) = 1 + k \cdot (1 - c)
\]

For $\cov(e,r) = 1$ (seen): boost $= 1$ (no change)\\
For $\cov(e,r) = 0$ (unseen): boost $= 1 + k$ (amplified)

The boost factor $k$ is \emph{learnable}:
\[
k = \exp(\phi)
\]
where $\phi$ is a scalar parameter. This allows the model to learn the optimal separation between seen and unseen pairs.

\paragraph{Triple Uncertainty.}
For a triple $(h, r, t)$, the total uncertainty is:
\[
U(h, r, t) = \sigma^2(h | r) + \sigma^2(t | r)
\]

This sums the relation-conditioned variances for head and tail entities.

\subsection{Training}

The loss function combines link prediction and uncertainty objectives:
\[
\mathcal{L} = \mathcal{L}_\text{score} + \lambda \cdot \mathcal{L}_\text{unc}
\]

\paragraph{Score Loss.}
Margin ranking loss for link prediction:
\[
\mathcal{L}_\text{score} = \sum_{(h,r,t) \in \mathcal{T}} \max(0, \gamma - s(h,r,t) + s(h,r,t'))
\]
where $t'$ is a corrupted (negative) tail and $\gamma$ is the margin.

\paragraph{Uncertainty Loss.}
Encourage higher uncertainty for negative samples:
\[
\mathcal{L}_\text{unc} = \sum_{(h,r,t) \in \mathcal{T}} \max(0, U(h,r,t) - U(h,r,t') + \epsilon)
\]

This trains the model to be more uncertain about negative (likely false) triples.

\subsection{Why Multiplicative Boost?}

We chose multiplicative over additive boost for theoretical and empirical reasons:

\paragraph{Scale Invariance.}
Multiplicative boost respects the scale of base variance:
\begin{itemize}
    \item High-variance entity in unseen context: even higher variance
    \item Low-variance entity in unseen context: moderately higher variance
\end{itemize}

Additive boost ($\sigma^2 + k$) would dominate low-variance entities while barely affecting high-variance ones.

\paragraph{Empirical Validation.}
Experiments confirm multiplicative outperforms additive by +2pp AUROC.

\subsection{Addressing the Circularity Concern}
\label{sec:circularity}

A natural concern is: ``If coverage defines what is OOD, then using coverage to detect OOD is circular---just use a lookup table.'' We address this directly:

\paragraph{Coverage-Only is Binary, Not Calibrated.}
A pure coverage detector $U(e,r) = 1 - \cov(e,r)$ achieves perfect AUROC (1.0) by construction. However, it provides \emph{no uncertainty gradations}: every OOD triple receives identical maximum uncertainty, and every ID triple receives zero uncertainty. This is a classifier, not an uncertainty estimator. RCUE combines perfect coverage-based detection with MLP-learned gradations, enabling meaningful uncertainty ranking \emph{within} the ID and OOD classes.

\paragraph{MLP Provides Semantic Calibration.}
The learned variance network $f_\theta$ captures how well-constrained an entity-relation combination is given embedding geometry. Within ID triples (where coverage provides no signal), the MLP achieves 0.66 AUROC for distinguishing ``easy'' vs ``hard'' predictions---16pp above Energy baselines. This semantic uncertainty is multiplicatively boosted for OOD samples, providing calibrated predictions across the full range.

\paragraph{Real-World Value Beyond Classification.}
In safety-critical deployments, binary detection is often insufficient. Systems need to \emph{rank} uncertainties: which OOD predictions are most likely wrong? RCUE provides this ranking via the MLP component, while coverage ensures all OOD samples receive elevated uncertainty. The coverage matrix is explicit engineering ($O(|\mathcal{E}| \cdot |\mathcal{R}|)$ hash table), not a learned capability---we are transparent about this.

\subsection{Computational Complexity}

RCUE adds minimal overhead:
\begin{itemize}
    \item \textbf{Parameters}: $O(2d \cdot h + h^2)$ for the MLP, plus one scalar $\phi$
    \item \textbf{Coverage matrix}: $O(|\mathcal{E}| \cdot |\mathcal{R}|)$ binary values
    \item \textbf{Inference}: One MLP forward pass per entity-relation pair
\end{itemize}

For large KGs, the coverage matrix can be replaced with a Bloom filter (5$\times$ compression, 0.14pp recall drop at 1\% FPR).
